\documentclass[main.tex]{subfiles}

\begin{document}

\newpage
\section{Tensor Products of Operator Systems}
\label{sec:tensor-products-of-operator-systems}

\begin{motivation}{\large Summary}{}
In this lecture, we will cover the following topics:
\begin{enumerate}
    \item The minimal, maximal, commuting, and injective tensor-product
    structures on operator systems.
    \item Nuclearity and the comparison problems that distinguish these
    tensor products.
\end{enumerate}
The principal source is Kavruk, Paulsen, Todorov, and Tomforde,
\emph{Tensor Products of Operator Systems} (\cite{kavruk-paulsen-todorov-tomforde-2011}). The contents of this lecture were given by Hridoyananda Saikia (\Cref{subsec:tensor_products_of_operator_systems_I}) and Lucas Hall (\Cref{subsec:tensor_products_of_operator_systems_II}). 
\end{motivation}

Here are the less common notational conventions used in this lecture.

\begin{filingbox}[NavyBlue]{Conventions}
If $S$ and $T$ are operator systems, then we denote the algebraic tensor product of their underlying vector spaces by $\tensor{\C}{S}{T}$.
\end{filingbox}

\newpage

Let's start with some motivation.

\begin{motiv}
Different physical and operator-algebraic models impose different positivity
conditions on the same algebraic tensor product.  Operator-system tensor
products make those choices explicit.  Their comparison maps and equality
problems are especially important in self-testing, where one must distinguish
tensor-product strategies from commuting-operator strategies without losing
the matrix-order structure seen by correlations.
\end{motiv}

\end{document}
